% docker run --rm -v $PWD:/workdir texlive/texlive:latest lualatex main.tex

% 印刷版ではPDF/X-4を生成する
% JapanColor2011Coated.icc はここから取得 → https://japancolor.jp/icc.html
% PDF/X-4と主張するファイルを出力するが、本当に準拠できてるかはまた別の問題
% 確認できないのでとりあえずコメントアウト
%\ifdefined\PRINT
%  \RequirePackage{pdfmanagement}
%  \SetKeys[document/metadata]{
%    lang=ja-JP,
%    pdfversion=1.6,
%    pdfstandard=X-4,
%    colorprofiles={X=JapanColor2011Coated.icc},
%  }
%\fi
% 最新の方法では \DocumentMetadata{pdfversion=1.6, pdfstandard=X-4,} と
% するようだが、jlreq との組み合わせではエラーになる

\input{revid}

% PDFに日付などの情報を埋め込まない = コンパイルするたびに同じPDFが出力される
% これだけでなく、hyperxmp パッケージも無効化するか、または環境変数SOURCE_DATE_EPOCHで
% 適当な時刻(epochからの経過秒)も合わせて指定する必要がある
% 32(CreationDate)+64(ModDate)+512(trailer ID)=608
%\pdfvariable suppressoptionalinfo 608

\ifdefined\PRINT
  \documentclass[book,
    fontsize=9pt, paper=a5paper, line_length=37zw, openany,
    twoside, gutter=22mm,% 印刷版: 両面、ノド側余白大きめ
  ]{jlreq}
\else
  \documentclass[book,
    fontsize=9pt, paper=a5paper, line_length=37zw, openany,
    oneside, % 電子版: 片面
  ]{jlreq}
\fi

\usepackage[haranoaji,       % 原の味フォントを使う
            no-math,         % 数式関連フォントはいじらない
            deluxe,          % 使えるウェイトを増やす
            expert,          % ルビなどで専用仮名フォントを使う
            match,           % \ttfamilyや\sffamilyで和文フォントをゴシックに
            jfm_yoko=jlreq,  % japanese font metric (よくわからんがおまじない)
]{luatexja-preset}
\usepackage[match]{luatexja-fontspec}
\ifdefined\PRINT
  \usepackage{jlreq-trimmarks}% トンボ
\fi
\usepackage{graphicx}        % includegraphics
\usepackage[most]{tcolorbox} % コラムなど
\usepackage{tikz}            % 描画
\usepackage{listings}        % ソースコード
\usepackage{tabularray}      % tblr環境
\usepackage{wallpaper}       % 表紙用
\usepackage{fontawesome7}    % アイコンフォント
\usepackage{qrcode}          % QRコード画像生成
\usepackage{eso-pic}         % 背景描画(表紙、隠しノンブル)
\usepackage{pageslts}        % 通しページ番号
% よく使いそうなパッケージ
%\usepackage{amsmath}        % align*
%\usepackage{tablefootnote}  % 表内で脚注(\tablefootnote)
%\usepackage{multicol}       % 指定箇所だけ段組
%\usepackage{wrapfig2}       % 図のまわりこみ
% PDF出力で \ref から \label にハイパーリンク/ \urlコマンド
% けっこうな数のコマンドを上書き定義しているので最後に読みこむのが推奨らしい
\usepackage[pdfusetitle]{hyperref}
\usepackage{hyperxmp}        % PDFメタデータ埋め込み/ hyperref のさらに後

\ifdefined\PRINT
  \selectcolormodel{gray}    % 1色刷り
  %\selectcolormodel{cmyk}   % カラー印刷
\else
  \selectcolormodel{rgb}
\fi

% \sffamilyをHelvetica系に(デフォのLatin Modern Sansは原の味角ゴと合わない)
\setsansfont{TeX Gyre Heros}
% 0/O、1/l/I が区別しやすく、高さ:幅=2:1の等幅フォント
\setmonofont{Inconsolatazi4}
% 原の味角ゴregularはInconsolataに合わせるとウェイトが重めなのでLightに
% ただ、Lightは軽すぎるような気もする。この中間ぐらいがほしい
\setmonojfont{HaranoAjiGothic-Light}

\ifdefined\PRINT
  % ページスタイル(jlreq専用) / 印刷版
  % 柱、ノンブルあり
  \NewPageStyle{body}{
    nombre_position=bottom-fore-edge, % ノンブル位置: 小口側下部
    nombre_font=\sffamily,
    running_head_position=top-fore-edge,% 柱位置: 小口側上部
    even_running_head={_chapter},  % 柱/左
    odd_running_head={_section},   % 柱/右
    mark_format={_chapter={第\thechapter 章\quad #1}, _section={\thesection\quad #1}}
  }
  % ノンブルのみ、柱なし
  \NewPageStyle{frontbackmatter}{
    nombre_position=bottom-fore-edge,
    nombre_font=\sffamily,
  }
\else
  % ページスタイル(jlreq専用) / 電子版
  % 柱、ノンブルあり
  \NewPageStyle{body}{
    nombre_position=bottom-center,% 下部中央
    nombre_font=\sffamily,
    odd_running_head={_chapter},% 柱: onesideではoddのみ有効っぽい
    mark_format={_chapter={第\thechapter 章\quad #1}}
  }
  % ノンブルのみ、柱なし
  \NewPageStyle{frontbackmatter}{
    nombre_position=bottom-center,
    nombre_font=\sffamily,
  }
\fi
% 本文ページスタイル(jlreqsetupで設定しても効かない？)
\pagestyle{body}

\jlreqsetup{
  frontmatter_pagestyle={frontbackmatter},
  mainmatter_pagebreak=clearpage, % \mainmatterで空白ページを挿入しない
  backmatter_pagestyle={frontbackmatter},
  % front/backmatter では章タイトル前後の余白を減らす/章番号を出力しない
  frontmatter_heading={ chapter={ before_lines=0, after_lines=0, number=false,},},
  mainmatter_heading={ chapter={ before_lines=*7, after_lines=3,},},
  backmatter_heading={ chapter={ before_lines=0, after_lines=0, number=false},},
}

% ページ通し番号(pagesLTS)
% \thepage は i,ii,iii,1,2,3,4,... となるので使わない
\pagenumbering{arabic}
% 文字列を1文字ずつ縦に並べる
% 要expl3; 現行lualatex/(u)platexなら動くはず(古いlatexエンジンでは動かないかも)
\ExplSyntaxOn
\NewDocumentCommand{\verticaltext}{O{1em}m}
  {
    \vbox{
      \hsize=#1
      \parindent=0pt
      \centering
      \tl_set:Nx \l_tmpa_tl {#2}
      \tl_map_inline:Nn \l_tmpa_tl
        {
          \hbox to #1{\hfil ##1\hfil}\par
        }
    }
  }
\ExplSyntaxOff
% 隠しノンブル
% ノド側(奇数ページは左、偶数ページは右）のページ端中央に配置
% 横組専用(縦組では左右が入れ替わるので位置調整が必要)
% 版面が用紙原点から1インチ内側に配置される前提(jlreq-trimmarksのデフォルト)
% jlreq-trimmmarksのオプションを変更したり、別の方法でトンボを出力したりするとズレる
% pagestyleとは独立しているので \thispagestyle{empty} でも出力される
% なんでこんな基本的なものを自前で実装せにゃならんのだ…
\ifdefined\PRINT
  \makeatletter
  \newcommand*{\ptlen}[1]{\strip@pt#1}
  \AddToShipoutPictureFG{%
    \begingroup
    \ttfamily\fontsize{6pt}{6pt}\selectfont
    \ifodd\value{page}% 奇数ページ(右ページ)
      \put(
        \ptlen\dimexpr 1in+1mm\relax,% ノド(左)から1mm内側
        \ptlen\dimexpr 1in+\paperheight/4\relax% 高さ中央(なぜ/2ではなく/4なのかわからん)
      ){% 背景に濃色の画像がある場合でもノンブルが読めるように数字の背景を白く塗り潰す
        {\setlength{\fboxsep}{0.5pt}\colorbox{white}{\verticaltext[3pt]{\theCurrentPage}}}%
      }%
    \else% 偶数ページ(左ページ)
      \put(
        \ptlen\dimexpr 1in+\paperwidth-1mm-1.4mm\relax,% ノド(右)から1mm+ノンブル幅(雑に1.4mmぐらい)内側
        \ptlen\dimexpr 1in+\paperheight/4\relax% 高さ中央
      ){%
        {\setlength{\fboxsep}{0.5pt}\colorbox{white}{\verticaltext[3pt]{\theCurrentPage}}}%
      }%
    \fi
    \endgroup
  }
  \makeatother
  % 裁ち切り線の外に出力されるバナーにもノンブルを追加
  % CurrentPage は実際のページ番号より1だけ進んだ値になっているので-1する
  \jlreqtrimmarkssetup{
    banner={
      top-left={
        yoko={page=\number\numexpr\value{CurrentPage}-1 \quad revid=\revid \quad jobname=\jobname},
      },
    },
  }
\fi

\definecolor{shadecolor}{gray}{0.9}  % snugshade環境その他

\definecolor{chaptergrey}{gray}{0.6} % 章番号
\newfontfamily{\chapnumfamily}{TeX Gyre Schola}
\newcommand{\chapnumfont}{\chapnumfamily\fontsize{100}{130}\selectfont\color{chaptergrey}\bfseries}

% 見出しの定義
% chapter; quotchap パッケージの模倣
\ModifyHeading{chapter}{
  pagestyle=frontbackmatter,% ノンブルのみ、柱なし
  font={},
  align=right,
  lines=7,         % 7行取りの中央に見出しを配置
  %before_lines=*7,% 見出し前後の空行(*7は改ページ直後でも強制的に7行空ける)
  %after_lines=3,  % \jlreqsetupで指定するのでコメントアウト
  after_label_space=0pt,
  second_heading_text_indent={*0pt},
  subtitle_indent={*0pt},
  format={% #1だと「第1章」になるので \thechapterで「1」だけ出力する
    \jlreqHeadingLabel{{\chapnumfont \thechapter\par\nobreak}}%
    {\huge\bfseries\raggedleft #2\par}%
  },%
}
% section; L字
\definecolor{sectionbar}{gray}{0.30}
\ModifyHeading{section}{
  font={\sffamily\bfseries\Large},
  label_format={\thesection},
  indent=0pt,% 行頭のアキ
  after_label_space=1em,% ラベルとタイトルの間のアキ
  second_heading_text_indent=0pt,
  lines=3,         % 見出しの高さ(ぜんぶで3行分の高さ)
  after_space=14pt,% 見出し後の空き
  format={%
    \noindent
    \begin{picture}(0,0)
      \put(0,-5){%
        \begin{tikzpicture}
          \fill[sectionbar] (0pt,0pt) rectangle (5pt,19pt);
        \end{tikzpicture}
      }
      \put(0,-5){%
        \color{sectionbar}
        \line(1,0){\dimexpr\hsize\relax}
      }
    \end{picture}%
    \hskip 1em
    #1#2%
  }%
 }
% subsection; 下線のみ
\ModifyHeading{subsection}{
  font={\large\sffamily\bfseries},
  label_format={\thesubsection},
  indent=0pt,
  after_label_space=1em,
  second_heading_text_indent=0pt,
  lines=2,
  after_space=8pt,
  format={%
    \noindent #1#2\par\nobreak
    \vskip2pt
    {\color{sectionbar}\hrule height 0.6pt depth0pt width\hsize}%
  }%
}
% subsubsection; 番号のかわりにアイコンを付加
\ModifyHeading{subsubsection}{
  number=false,
  font={\sffamily\bfseries},
  indent=0pt,
  lines=1,
  second_heading_text_indent=0pt,
  format={\noindent\faIcon{book}\quad #2}%
}

% 図表見出し
\renewcommand{\tablename}{\textcolor{gray}{▼} 表}
\renewcommand{\figurename}{\textcolor{gray}{▲} 図}
\renewcommand{\lstlistingname}{\textcolor{gray}{▼} リスト}

% ソースコード
\lstset{
%  language = bash,		% 言語
  backgroundcolor={\color[gray]{.90}},	%背景色と透過度
  breaklines = true,		%枠外に行った時の自動改行
  breakindent = 0pt,		%自動改行後のインデント量(デフォルトでは20[pt])	
  %basicstyle = \ttfamily\footnotesize,		% 標準の書体
  basicstyle = \ttfamily\small,		% 標準の書体
  %commentstyle = {\itshape \color[cmyk]{1,0.4,1,0}},	%コメントの書体
  %commentstyle = \ttfamily,
  classoffset = 0,		%関数名等の色の設定
  %keywordstyle = {\bfseries \color[cmyk]{0,1,0,0}},	%キーワード(int, ifなど)の書体
  %keywordstyle = {\bfseries},
  stringstyle = \ttfamily,
  showstringspaces=false,
  frame = tBRl,	%枠 "t"は上に線を記載, "T"は上に二重線を記載
		%他オプション：leftline，topline，bottomline，lines，single，shadowbox
  framesep = 5pt,	%枠と内部の文字の間隔
  numbers = left,	%行番号の位置
  numbersep = 9pt,     %枠と行番号の間隔
  stepnumber = 1,	%行番号の間隔
  numberstyle = \tiny,	%行番号の書体
  tabsize = 8,		%タブの大きさ
  basewidth = 0.5em,	%字送り
  captionpos = t	%キャプションの場所("tb"ならば上下両方に記載)
}

% 実行結果表示用(上の設定からの差分のみ)
\lstdefinestyle{tty}{
  basicstyle = \ttfamily\footnotesize,
  frame= {},
  framesep = 3pt,
  numbersep = 5pt,
}

% コラム(tcolorbox)
\definecolor{columnframe}{gray}{0.40}
\DeclareTColorBox{column}{m}{
  title={\faIcon{comment-dots}\quad #1},
  fonttitle=\sffamily,
  breakable, % コラム内で改ページできるようにする
  before upper={ % 本文直前に挿入されるコード
    \setlength{\parindent}{1\zw}% 最初の段落で字下げする(デフォではされない)
  },
  before={ % 箱の直前に挿入されるコード
    \vspace{\baselineskip}
    \phantomsection% hyperrefの被リンク用ダミーセクションを作る
    \addcontentsline{toc}{section}{\numberline{\faIcon[regular]{comment-dots}}#1}%目次に載せる
  },
  colframe=columnframe, % 上の方で\definecolor した色
  boxrule=0.5pt,
  arc=1.5mm,
  enhanced, % underlayでいろいろできるようにする
  underlay unbroken={ % 1ページに収まる場合
    \draw[line width=0.5pt, rounded corners=1.5mm]
      (frame.north west) rectangle (frame.south east);
  },
  underlay first={ % 分割された最初のパーツ
    \draw[line width=0.5pt, rounded corners=1.5mm]
      (frame.south west) -- (frame.north west) -- (frame.north east) -- (frame.south east);
  },
  underlay middle={ % 分割された途中のパーツ
    \draw[line width=0.5pt] (frame.north west) -- (frame.south west);
    \draw[line width=0.5pt] (frame.north east) -- (frame.south east);
  },
  underlay last={ % 分割された最後のパーツ
    \draw[line width=0.5pt, rounded corners=1.5mm]
      (frame.north west) -- (frame.south west) -- (frame.south east) -- (frame.north east);
  },
}

% framed package に含まれる snugshade, leftbar 環境に似せたもの(完全には同じではない)
\DeclareTColorBox{snugshade}{}{
  breakable,
  before upper={
    \setlength{\parindent}{1\zw}
  },
  width=\linewidth,
  top=2pt,bottom=2pt,left=2pt,right=2pt,
  colframe=shadecolor,
  colback=shadecolor,
  boxrule=0pt, arc=0pt,
}
\DeclareTColorBox{leftbar}{}{
  breakable,
  before upper={
    \setlength{\parindent}{1\zw}
  },
  width=\linewidth,
  top=2pt,bottom=2pt,left=8pt,right=2pt,
  colframe=white,
  colback=white,
  boxrule=0pt, arc=0pt,
  enhanced,
  underlay={
    \draw[black, line width=2pt](frame.north west) -- (frame.south west);
  },
}

% tcolorbox の中で \footnote を使うと、番号が1),2),3),...ではなくa,b,c,...になる
% これを *),  \dagger), \ddagger),... になるように変更
% mpfootnote は minipage 環境用のカウンタらしい
\renewcommand{\thempfootnote}{\fnsymbol{mpfootnote})}

\title{かっこいいタイトル}
\author{何野誰某}
\date{YYYY/MM/DD}

% その他の hyperref 設定
\hypersetup{
% luatex,% 自動検出されるので不要
% unicode=true,% これもデフォなので不要
% pdfusetitle,% \title, \author を埋め込む(\usepackageでの指定が必須なのでここではコメントアウト)
  bookmarks=true, % 見出し
  bookmarksdepth=2,  % デフォではtocdepthと同じ値だけどいちおう
  hidelinks, % リンクに色・枠をつけない
%  pdfkeywords={key, word},
%  pdfsubject={pdf subject},% titleとは別っぽい
  pdflang=ja-JP,
  % 以下、hyperxmp
%  pdfcontactemail={mail@address},
%  pdfcontacturl={https://example.com/},
  pdfversionid=\revid,
%  pdfdate={1970-01-01},
%  keeppdfinfo,% メタ情報を dc:hoge だけでなくInfo dictにも出力する
}

\begin{document}
\frontmatter
\ifdefined\PRINT
%  % 扉画像ページ(印刷版)
%  % 印刷版PDFには表紙画像を含めず、本文とは別に入稿するのが基本
%  % 表紙をめくったら1ページ目に扉としてモノクロ版の画像を入れるような場合には
%  % コメントアウトされている部分を生かす
%  \thispagestyle{empty} % emptyでも隠しノンブルは入る
%  % トンボ(1in)と塗り足し幅(3mm)のぶんだけ位置がズレるので補正
%  % 1inという値はjlreq-trimmarksのデフォルトが2in大きい用紙を使うことに由来
%  % デフォルトから変更していたり、別の方法でトンボを作る場合には要修正
%  \setlength{\wpXoffset}{22.4mm}  % 1in-3mm
%  \setlength{\wpYoffset}{-28.4mm} % -1in-3mm
%  % 画像は版面いっぱいではなく、塗り足し幅(3mm)の分拡大する
%  % (A5幅148mm + 塗り足し幅3mm x 2) / 148mm = 1.041; B5なら(182+6)/182=1.033
%  % 画像ファイルは縦横比1:1.414ではなく、154:216 (B5なら188:263)の比率で
%  % グレースケール(またはCMYK)の高解像度のものを指定する
%  % A5グレースケール600dpiドットバイドットなら3638x5102ピクセル
%  \ThisCenterWallPaper{1.041}{cover/print.png}
%  \null\clearpage
\else
  % 表紙(電子版)
  % \ThisCenterWallPaper の最初の引数は版面の横幅に対する画像幅の比率
  % 1なら版面にちょうどぴったり収まるよう画像サイズを調整する
  % 縦横比は保持されるので、1:1.414 の比率の画像ファイルを指定すると
  % 縦もぴったり版面に収まる
  %\setcounter{page}{0} % 表紙画像をページ数に数えない場合
  \thispagestyle{empty}
  \ThisCenterWallPaper{1}{cover/cover.png}
  \null% この\nullがないと表示されない
  \clearpage
\fi

% 序文
\ifdefined\PRINT
%  % 扉画像ページ(印刷版)
%  % 印刷版PDFには表紙画像を含めず、本文とは別に入稿するのが基本
%  % 表紙をめくったら1ページ目に扉としてモノクロ版の画像を入れるような場合には
%  % コメントアウトされている部分を生かす
%  \thispagestyle{empty} % emptyでも隠しノンブルは入る
%  % トンボ(1in)と塗り足し幅(3mm)のぶんだけ位置がズレるので補正
%  % 1inという値はjlreq-trimmarksのデフォルトが2in大きい用紙を使うことに由来
%  % デフォルトから変更していたり、別の方法でトンボを作る場合には要修正
%  \setlength{\wpXoffset}{22.4mm}  % 1in-3mm
%  \setlength{\wpYoffset}{-28.4mm} % -1in-3mm
%  % 画像は版面いっぱいではなく、塗り足し幅(3mm)の分拡大する
%  % (A5幅148mm + 塗り足し幅3mm x 2) / 148mm = 1.041; B5なら(182+6)/182=1.033
%  % 画像ファイルは縦横比1:1.414ではなく、154:216 (B5なら188:263)の比率で
%  % グレースケール(またはCMYK)の高解像度のものを指定する
%  % A5グレースケール600dpiドットバイドットなら3638x5102ピクセル
%  \ThisCenterWallPaper{1.041}{cover/print.png}
%  \null\clearpage
\else
  % 表紙(電子版)
  % \ThisCenterWallPaper の最初の引数は版面の横幅に対する画像幅の比率
  % 1なら版面にちょうどぴったり収まるよう画像サイズを調整する
  % 縦横比は保持されるので、1:1.414 の比率の画像ファイルを指定すると
  % 縦もぴったり版面に収まる
  %\setcounter{page}{0} % 表紙画像をページ数に数えない場合
  \thispagestyle{empty}
  \ThisCenterWallPaper{1}{cover/cover.png}
  \null% この\nullがないと表示されない
  \clearpage
\fi

% 序文
\chapter{序}

なんかかっこいいことを書く。


% 目次(subsectionまで)
\setcounter{tocdepth}{2}
\tableofcontents

% 本文
\mainmatter
\chapter{README}

章タイトルのデザインは quotchap パッケージ\footnote{\url{https://ctan.org/pkg/quotchap}}のパクリ。

\section{執筆のしかた}

このリポジトリを \verb`git clone` したら、\verb`*.tex` をてきとーに編集して、\verb`make` するだけ。

必要に応じて、\verb`git remote remove origin` してから自前のリポジトリに \verb`git remote add origin ...` しなおすべし。

\subsection{\LaTeX 環境}

Lua\LaTeX + jlreq が必要。\TeX Live で動くはず。

docker で構築するのが楽ちん。

\begin{lstlisting}[style=tty]
% docker pull texlive/texlive:latest
\end{lstlisting}

docker を使わない場合は \verb`Makefile` の \verb`LATEXMK=...` の行を修正する必要がある。

\begin{lstlisting}
# docker を使う場合
LATEXMK= docker run --rm -v .:/workdir texlive/texlive:latest latexmk
# 使わない場合
#LATEXMK= latexmk
# ローカルではdockerを使わないが、github actions でコンパイルするときには使う
#LATEXMK= $(shell command -v latexmk || echo "docker run --rm -v .:/workdir texlive/texlive:latest latexmk")
\end{lstlisting}

\subsection{コンパイル}

\verb`make` 一発。

\verb`Makefile` をとくに修正していない場合、電子版(画面上での閲覧を想定したPDF)の \verb`main.pdf` を生成する。また、\verb`make all` ではそれに加えて印刷版(印刷所への入稿を想定したPDF)の \verb`print.pdf` を生成する。コンパイルのたびに両方生成するのは時間のムダなので、デフォルトではあえて片方だけ生成するようにしている。

\verb`make` の際、\verb`git` のコミットハッシュが記載された \verb`revid.tex` というファイルを自動で生成する。この値は \verb`\revid` というマクロで参照できる。本ドキュメントでは奥付で使用している。また、PDFのメタ情報にも埋め込んでいる。

\begin{lstlisting}[style=tty]
% git rev-parse --short main
1eb445e
% cat revid.tex
\def\revid{1eb445e}
% pdfinfo -meta main.pdf | fgrep VersionID
                  <pdfaProperty:name>VersionID</pdfaProperty:name>
      <xmpMM:VersionID>1eb445e</xmpMM:VersionID>
\end{lstlisting}

画像を読み込むための \verb`\includegraphics{}` とソースコードを読み込むための \verb`\lstinputlisting{}` で指定されているファイルは自動的に依存関係に追加されるので(\verb`Makefile` の \verb`INCLUDE=...` の行)、これらのファイルを修正すれば \verb`*.tex` を修正していなくても \verb`make` で再生成される。ただし、環境変数 \verb`$TEXINPUTS` で指定されたディレクトリまでは見ない。\verb`$TEXINPUTS` に頼らない書き方をするのがよい。

\verb`make` が使えない環境(Windows や Overleaf のような Web 上のもの)は知らん。がんばって何とかしてくれ。

\subsection{github workflow}

\verb`main` または \verb`master` ブランチで github に push すると、github 上で \verb`make all` される。docker を使わずにコンパイルするように \verb`Makefile` の \verb`LATEXMK=...` の行を修正していた場合、github でのビルドがコケるので注意。

自動コンパイルしたくない場合は \verb`.github/` 以下をまるごと削除すべし。

\subsection{印刷版と電子版}

前述のように \verb`make all` で2種類のPDFを生成する。主な違いは以下。

\begin{itemize}
\item 印刷版は見開き左右でページデザインが異なるが(twoside)、左右の概念がない電子版はどのページも同じデザイン(oneside)
\item 電子版は1ページ目に表紙画像を挿入するが、印刷版では別ファイルでの入稿を想定して挿入しない
\item 印刷版はグレースケール\footnote{ただし、カラー画像がグレースケールに変換されるわけではない}、電子版はRGBカラー
\item 印刷版ではトンボと隠しノンブルを出力する
\end{itemize}

これ以外に電子版、印刷版で異なる処理をさせたい場合は、\verb`main.tex` で以下のように記述する。

\begin{lstlisting}
\ifdefined\PRINT
  % 印刷版の処理
\else
  % 電子版の処理
\fi
\end{lstlisting}

これは\ifdefined\PRINT 印刷版\else 電子版\fi のPDFである。

\verb`print.tex` は電子版か印刷版かを区別するための \verb`\PRINT` マクロを定義してから \verb`main.tex` を読み込むだけのファイルで、通常はとくに修正する必要はない。

\section{セクション}

ここからは見た目を確認するためのデモ。

\subsection{サブセクション}

びよーん。

\subsubsection{サブサブセクション}

ぱおーん。

\begin{column}{コラムタイトル}

コラムはこう書く\footnote{コラム内で脚注を書くとこうなる}。

目次にも自動的に掲載される。

\end{column}

ソースコードはリスト\ref{src}のように書く。

\begin{lstlisting}[caption=hello.sh, label=src]
echo Hello, world!
\end{lstlisting}

枠のないバージョン。

\begin{lstlisting}[style=tty]
$ sh hello.sh
Hello, world!
\end{lstlisting}

なんか引用文を入れるのによさそうな環境。えらい人曰く、

\begin{leftbar}
腹へった。
\end{leftbar}

\subsection{フォント}

\def\sampletext{\noindent%
A quick brown fox jumps over the lazy dog. \\
色は匂へど散りぬるを我が世誰ぞ常ならむ有為の奥山今日越えて浅き夢見し酔ひもせず \\
0123456789
}

\textrm{\textbackslash{}rmfamily}
\begin{snugshade}
\rmfamily\sampletext
\end{snugshade}

\textsf{\textbackslash{}sffamily}
\begin{snugshade}
\sffamily\sampletext
\end{snugshade}

\texttt{\textbackslash{}ttfamily}
\begin{snugshade}
\ttfamily\sampletext
\end{snugshade}



% あとがき・奥付
\backmatter
\chapter{あとがき}

本文と関係ないことで無駄に1ページを費すことに対する反省の弁。


\clearpage
\thispagestyle{empty}
% ここから下は、かなり雑にレイアウトしてるので適宜修正すべし。

\begin{flushleft}
{\LARGE\faIcon{pen-to-square}}\quad{\large\sffamily 著者紹介}

\vspace{0.5\baselineskip}
{\sffamily{}何野誰某}
\vspace{0.3\baselineskip}

そのへんのおっさん。
\end{flushleft}

\vspace{1\baselineskip}% QRコードの周囲にほとんど空白が空かないので手で調整する
\begin{tblr}{cc}
\qrcode[height=2cm]{https://x.com/example_account} &
\qrcode[height=2cm]{https://example.com} \\
\href{https://x.com/example_account}{\faIcon{square-x-twitter} \ttfamily@example\_account} &
\url{https://example.com}
\end{tblr}
\vspace{1\baselineskip}



% 奥付は偶数ページに載せる
% 以下のコメント部分を生かすと自動で判定されるが、手作業で調整してもいいんじゃないかな
% (偶数ページなら改ページせず著者紹介の下のスペースに、奇数ページなら改ページ)
%\makeatletter
%\ifodd\c@CurrentPage
%  \clearpage\thispagestyle{empty}
%\fi
%\makeatother

\null\vfill

% 奥付本体
\begin{center}
\parbox{25\zw}{
{\large\sffamily{}かっこいいタイトル}
\vspace*{0.5\zh}
\hrule height 1pt % 太さ1pt横幅いっぱいの線
\vspace*{1\zh}
YYYY年MM月DD日 初版第1刷
\vspace*{1\zh}
\begin{center}
\begin{tblr}{ll}
著者 & 何野誰某 \\
発行 & Foobar Publisher \\
e-mail & \texttt{mail@address} \\
印刷 & ほげほげ印刷 \\
\end{tblr}
\end{center}
\vspace*{1\zh}
\hrule height 0.4pt
\null\hfill{\ttfamily rev. \revid}
}% end \parbox
\end{center}


% 裏表紙
%\ifdefined\PRINT
%  % 印刷版では裏表紙は別ファイルとして入稿
%\else
%  % 電子版
%  \thispagestyle{empty}
%  \ThisCenterWallPaper{1}{cover/cover4.png}
%  \null
%\fi

\end{document}
